\medskip

On considère le cube ABCDEFGH de côté 1 représenté ci-dessous.

\begin{center}
\psset{unit=1cm}
\begin{pspicture}(7.2,6.9)
\psframe(0.2,0.2)(5.2,5.2)
\psline(5.2,0.2)(6.8,1.5)(6.8,6.5)(5.2,5.2)
\psline(6.8,6.5)(1.8,6.5)(0.2,5.2)
\psline[linestyle=dashed](0.2,0.2)(1.8,1.5)(6.8,1.5)
\psline[linestyle=dashed](1.8,1.5)(1.8,6.5)
\uput[dl](0.2,0.2){A} \uput[dr](5.2,0.2){B} \uput[r](6.8,1.5){C} \uput[l](1.8,1.5){D}
\uput[l](0.2,5.2){E} \uput[r](5.2,5.2){F} \uput[ur](6.8,6.5){G} \uput[ul](1.8,6.5){H}
\end{pspicture}
\end{center}

On munit l'espace du repère orthonormé $\left(\text{A}~;~\vect{\text{AB}},\, \vect{\text{AD}},\, \vect{\text{AE}}\right)$.

\medskip

\begin{enumerate}
\item 
	\begin{enumerate}
		\item Justifier que les droites (AH) et (ED) sont perpendiculaires.
		\item Justifier que la droite (GH) est orthogonale au plan (EDH).
		\item En déduire que la droite (ED) est orthogonale au plan (AGH).
	\end{enumerate}	
\item Donner les coordonnées du vecteur $\vect{\text{ED}}$.

Déduire de la question 1. c. qu'une équation cartésienne du plan (AGH) est:

\[y - z = 0.\]

\item On désigne par L le point de coordonnées $\left(\dfrac23~;~1~;~0\right)$.
	\begin{enumerate}
		\item Déterminer une représentation paramétrique de la droite (EL).
		\item Déterminer l'intersection de la droite (EL) et du plan (AGH).
		\item Démontrer que le projeté orthogonal du point L sur le plan (AGH) est le
point K de coordonnées $\left(\dfrac23~;~\dfrac12~;~\dfrac12\right)$.
		\item Montrer que la distance du point L au plan (AGH) est égale à $\dfrac{\sqrt{2}}{2}$.
		\item Déterminer le volume du tétraèdre LAGH.
		
On rappelle que le volume $V$ d'un tétraèdre est donné par la formule :

\[V = \dfrac13 \times (\text{aire de la base}) \times \text{hauteur}.\]
	\end{enumerate}
\end{enumerate}

\bigskip
